\documentclass{article}
\usepackage{amsmath}
\usepackage{amssymb}
\usepackage{booktabs}
\usepackage{array}
\usepackage[a4paper,margin=1in]{geometry}

\newcolumntype{L}[1]{>{\raggedright\arraybackslash}p{#1}}

\title{Choosing the Promoter-Mining Ensemble\\ \large How six DNA models were compared, and why three were deployed}
\date{September 2026}
\author{HKUST iGEM 2026 Dry Lab}

\begin{document}
\maketitle

\begin{abstract}
We evaluated six machine-learning models for finding bacterial promoters in the genome of our target,
\emph{Clostridium butyricum}, and chose the best three for the deployed ensemble. The headline result is
that a model's score on a standard \emph{balanced} test is a poor guide to how it behaves at \emph{genome
scale} --- the setting that actually matters when mining a real genome. On two held-out \emph{Clostridium}
genomes, a DNA foundation-model probe (Evo2), a fine-tuned language model (ProkBERT), and a custom fusion
CNN were the three most precise, and combining them into one ensemble beat every individual model. Two
Nucleotide Transformer variants and a simpler CNN were evaluated and set aside.
\end{abstract}

\section{What we measured}
Six models were each tested two ways. The first is a \textbf{balanced test}: equal numbers of real
promoters and non-promoters, the standard machine-learning benchmark (scored by the area under the
precision--recall curve). The second is a \textbf{genome-scale test}, which slides each model across an
entire genome held out of training, where real promoters make up only a tiny fraction of all positions
--- the realistic mining setting.

For the genome-scale test we report \textbf{precision at a fixed recall}. \emph{Recall} is the fraction of
real promoters the model finds; \emph{precision} is the fraction of its predictions that are correct. So
``precision at 50\% recall'' reads: when the model is tuned to recover half of the real promoters, what
fraction of the sites it flags are genuine? Higher is better --- it means fewer false alarms. We use two
\emph{Clostridium} genomes held out of training, \emph{C.\ ljungdahlii} and \emph{C.\ autoethanogenum},
as close-relative stand-ins for the target.

The six models are a DNA foundation-model probe (\textbf{Evo2}: a lightweight classifier on top of the
Evo2 model from the Arc Institute, which is used as-is, not retrained), a fine-tuned prokaryotic language
model (\textbf{ProkBERT}), a custom convolutional network trained from scratch that fuses several views of
the sequence (\textbf{Fusion CNN}), a plainer from-scratch convolutional network as a control
(\textbf{Simple CNN}), a DNA transformer fine-tuned with LoRA (\textbf{Nucleotide Transformer}), and a
larger DNA transformer that reads a wider window (\textbf{Nucleotide Transformer v3}).

\section{Balanced accuracy is not genome-scale performance}
On the balanced test the ranking is almost the \emph{reverse} of what matters at genome scale. Only one
model clears 0.90 --- the Nucleotide Transformer v3, at 0.96 (Table~\ref{tab:balanced})\footnote{Not
directly comparable to the others: v3 reads a wider input window than the 81-base-pair models, which
inflates its balanced score.} --- and it is the \emph{weakest} of the six at genome scale
(Table~\ref{tab:genome}). The Evo2 probe is its mirror image: \emph{lowest} of all on the balanced test
(0.70), yet the \emph{strongest} at genome scale. This reversal is the central lesson of the evaluation,
and the reason we chose the deployment models on genome-scale behaviour rather than balanced accuracy.

\begin{table}[h]
\centering
\caption{Balanced-test score (area under the precision--recall curve; higher is better). Only the
Nucleotide Transformer v3 clears 0.90 --- and Table~\ref{tab:genome} shows the genome-scale ranking is
nearly the reverse.}
\label{tab:balanced}
\begin{tabular}{lc}
\toprule
Model & Balanced-test score \\
\midrule
Nucleotide Transformer v3 & 0.960 \\
ProkBERT                  & 0.884 \\
Nucleotide Transformer    & 0.854 \\
Fusion CNN                & 0.848 \\
Simple CNN                & 0.804 \\
Evo2 probe                & 0.703 \\
\bottomrule
\end{tabular}
\end{table}

\begin{table}[h]
\centering
\caption{Genome-scale precision on two held-out \emph{Clostridium} genomes, at 10\% and 50\% recall.
Higher is better (fewer false positives). \checkmark\ marks the three models chosen for the deployed
ensemble.}
\label{tab:genome}
\begin{tabular}{lccccc}
\toprule
& & \multicolumn{2}{c}{\emph{C.\ ljungdahlii}} & \multicolumn{2}{c}{\emph{C.\ autoethanogenum}} \\
\cmidrule(lr){3-4}\cmidrule(lr){5-6}
Model & Deployed & @10\% & @50\% & @10\% & @50\% \\
\midrule
Evo2 probe                & \checkmark & \textbf{71.0\%} & \textbf{58.0\%} & \textbf{85.8\%} & \textbf{71.4\%} \\
ProkBERT                  & \checkmark & 64.4\% & 47.7\% & 76.7\% & 64.6\% \\
Fusion CNN                & \checkmark & 64.4\% & 47.2\% & 81.0\% & 66.0\% \\
Simple CNN                & ---        & 49.0\% & 40.3\% & 72.6\% & 63.0\% \\
Nucleotide Transformer    & ---        & 43.9\% & 36.1\% & 64.9\% & 58.7\% \\
Nucleotide Transformer v3 & ---        & 35.7\% & 28.0\% & 58.5\% & 48.3\% \\
\bottomrule
\end{tabular}
\end{table}

The genome-scale ranking is stable across both genomes: Evo2 leads, ProkBERT and the Fusion CNN follow
closely, and the two Nucleotide Transformers and the simple CNN trail --- the exact opposite of the
balanced-test order, where Evo2 is last.

\section{The ensemble beats every single model}
Rather than pick one model, we combine the top three by agreement: each ranks every candidate site, and we
average those ranks. Because it depends only on the ordering of sites, this combination is unaffected by
the models scoring on different scales. The result is the most precise caller on both held-out genomes.

\begin{table}[h]
\centering
\caption{The three-model ensemble versus its members: genome-scale precision at 10\% and 50\% recall.}
\label{tab:ensemble}
\begin{tabular}{lcccc}
\toprule
& \multicolumn{2}{c}{\emph{C.\ ljungdahlii}} & \multicolumn{2}{c}{\emph{C.\ autoethanogenum}} \\
\cmidrule(lr){2-3}\cmidrule(lr){4-5}
& @10\% & @50\% & @10\% & @50\% \\
\midrule
\textbf{Three-model ensemble} & \textbf{74.4\%} & \textbf{61.6\%} & \textbf{87.0\%} & \textbf{76.6\%} \\
Evo2 probe                    & 69.4\% & 58.4\% & 87.6\% & 72.5\% \\
ProkBERT                      & 64.4\% & 47.7\% & 76.7\% & 64.6\% \\
Fusion CNN                    & 64.4\% & 47.2\% & 81.0\% & 66.0\% \\
\bottomrule
\end{tabular}
\end{table}

The ensemble matches or beats the best single model at every recall level, and is decisively ahead at
50\% recall on both genomes. Adding a fourth model made it \emph{worse} on both (for example, precision at
50\% recall on \emph{C.\ ljungdahlii} fell from 61.6\% to 58.3\%), confirming that three complementary
models is the right number.\footnote{Single-model figures differ by about a percentage point between
Tables~\ref{tab:genome} and~\ref{tab:ensemble} because they come from two separate evaluation runs; the
comparisons \emph{within} each table are the valid ones.}

\section{What we deployed}
The final ensemble is three deliberately different models --- one DNA foundation model (Evo2), one
language model (ProkBERT), and one convolutional network (the Fusion CNN) --- so their independent
strengths reinforce one another. This ensemble was then used to mine the target genome, \emph{C.\
butyricum}, producing a ranked shortlist of candidate promoters for experimental validation.

\end{document}
